% =====================================================================
%  Section: Masses   (Plan Stages 5-7)
%  Paper 1, "It from Bit via Gödel". Drafted 2026-06-08.
%  Follows \S sec:generations; precedes the boson-mass / electroweak section.
%
%  \input fragment. biblatex/\autocite, amsthm `proposition'.
%  New references: koide1983new, georgijarlskog1979new
%  (delivered in masses_refs_2026-06-08.ris). J3(O) via baez2002octonions
%  and duboisviolettetodorov2019exceptional, already in the library.
%
%  Cross-reference labels (placeholders -- re-point before compiling):
%    sec:generations  -- the particle zoo / depth ordering (previous section)
%    sec:charges      -- quantum numbers; the "debt-free" photon forward-ref lands here
%    sec:register     -- the three-qubit register
%    sec:interactions -- circuits branch (forward; CKM, FCNC)
%    sec:boson-masses -- boson masses / electroweak parameters (forward)
%    sec:predictions  -- neutrino masses, glueballs (forward)
%
%  HONESTY FLAGS, now recorded as named open problems (paper-wide
%  convention; preamble in the quantum-numbers header):
%    op:associator-geodesic -- associator-debt vs geodesic proxy;
%    op:koide-delta         -- derive delta = 2/9 from Hopf holonomy
%                              (key open calculation; sin^4 holonomy route
%                              recorded as tried-and-failed);
%    op:colour-correction   -- derive |Q|^{3/2} law + phase-dilution rule;
%                              reconcile Casimir-vs-fitted Georgi-Jarlskog;
%    op:up-quark            -- the +29% discrepancy (candidate mechanism:
%                              generation-1 GHZ/W interference);
%    op:sector-scales       -- reduce the three sector scales to one;
%    op:cabibbo             -- mass phase = Cabibbo angle.
%
%  REVISED 2026-06-11 after review and full numerical re-verification:
%  lepton table corrected to source values (e -0.01%, mu -0.01%; drafted
%  +0.2%/-0.03% were stale); PHASE-DILUTION RULE delta_sector = (2/9)/n
%  added to the quark subsection -- its omission made the quoted quark
%  numbers unreproducible (naive same-delta transplant fails by orders of
%  magnitude); 2/9 dimension reading corrected: dim(OP^1) = 8, not 9; the
%  correct nine-dimensional object is O+R, the ambient of the Hopf base
%  (= Bloch + colour, app:lem:base). BACKPORT: source docs carry
%  dim(OP^1) = 9 -- fix there too. One-scale claim made honest (three
%  sector scales as computed; y_t = 1 ties up-sector to v). Neutrino
%  forward-pointer: RESOLVED (2026-06-12) to ~52 meV per
%  fermion_mass_geodesic_calculation.md (dilution rule, n = 2); the
%  2.6 meV figure in quark_neutrino_masses.md is the superseded early fit
%  -- BACKPORT a supersession note there. sec:predictions (to be drafted)
%  carries: Sigma ~52 meV with the Delta-m^2_21 (x30) tension and the
%  0.6-degree theta_nu fine-tuning as a named open problem; PLUS the
%  robust qualitative set, stated as such: m_1 = 0 exactly (rank-2 seesaw
%  from two steriles), normal ordering mandatory, two sterile states
%  ~3.6 TeV, no fourth generation at any scale; falsifiers JUNO, DUNE,
%  DESI, Euclid, CMB-S4.
% =====================================================================

\section{Masses}
\label{sec:masses}

The previous section enumerated the spectrum and asserted, qualitatively, that mass grows with the
algebraic depth of the frustrated qubit. This section makes that quantitative. It begins by saying
what mass \emph{is} in this framework --- the promised definition of the ``debt'' whose absence made
the photon massless in \S\ref{sec:charges} --- and then computes two spectra: the charged leptons,
which are the clean showcase, reproduced from a single dimensional scale and no free dimensionless
parameters; and the quarks, which require two structural modifications --- one to the amplitude, one
to the phase --- and succeed for eight of the nine charged fermions, with the one failure flagged
plainly. Throughout, masses are computed on the canonical orbits of \S\ref{sec:canonical-orbits},
where the reduced spectra --- and with them the quantities below --- are sharp.

\subsection{What mass is: associator-debt}

In an associative algebra the product $(ab)c$ equals $a(bc)$ and there is nothing to record. The
octonions are not associative, and the failure is measured by the associator
\begin{equation}
  [a,b,c] = (ab)c - a(bc),
\end{equation}
which is generically non-zero. The framework identifies \emph{mass with this associator-debt}: the
information cost of specifying a bracketing that the algebra itself leaves ambiguous. Dimensionally
the associator carries an inverse length in natural units, as a mass must. A channel that carries no
associator-debt is massless --- this is exactly the ``debt-free'' direction of \S\ref{sec:charges},
and it is why the photon, the one electroweak direction with vanishing debt, has no mass.

It is important to be exact about what the associator does and does not fix. Computed on the basis
units, the associator has a uniform magnitude across all bracketings that are not forced to vanish;
it therefore establishes that mass exists and is non-zero, and that the three generations'
bracketings are inequivalent, but it does \emph{not} by itself set the mass hierarchy. The hierarchy
comes from how the triality structure breaks --- the angle of the next subsection --- while the
associator supplies the fact of mass. This division of labour, associator for existence and triality
angle for ratios, is the honest reading of the calculation and we keep the two separate throughout.

Two structural points complete the picture. First, because the non-associativity forbids a metric on
the octonionic fibre $S^7$, mass cannot literally be a geodesic length on that fibre; it is the
associator-debt. The operational calculations below, which are sometimes phrased as ``geodesic''
quantities, are a computational proxy.

\begin{openproblem}[Associator-debt and the geodesic proxy]\label{op:associator-geodesic}
Show that the geodesic-length proxy used in the operational mass calculations computes the same
quantity as the associator-debt that defines mass, closing the gap between definition and
computation.
\end{openproblem}

Second, the dimensional-parameter count. Structurally the framework requires exactly one free
dimensional parameter --- an overall scale --- and this is a necessity rather than a gap: the
observer cannot fix the absolute scale of their own computational network from inside it
(\S\ref{sec:register}). The Higgs vacuum value sets that scale, and the doubling norm $\sqrt2$ that
appears below is the algebraic signature of the vacuum value ``doubling'' the structure. As the
calculation currently stands, however, each charge sector carries its own fitted scale --- three in
all --- of which the up-quark scale is tied to the electroweak value by $y_t = 1$
(\S\ref{sec:boson-masses}), while the ratio of the down-quark to the lepton scale sits suggestively
near the doubling norm squared, $M^2_{\mathrm{down}}/M^2_{\mathrm{lepton}} \approx 2.07$. Reducing
the three to the structurally required one is open.

\begin{openproblem}[The sector scales]\label{op:sector-scales}
Derive the lepton and down-quark sector mass scales from the single electroweak scale, as
$y_t = 1$ already does for the up-quark sector --- or show why the residual ratios
($\approx 2$ for down-to-lepton) are fixed by the algebra.
\end{openproblem}

\subsection{Charged leptons: the Koide relation}

The charged-lepton masses satisfy the empirical relation of Koide \autocite{koide1983new},
\begin{equation}
  \frac{m_e + m_\mu + m_\tau}{\bigl(\sqrt{m_e}+\sqrt{m_\mu}+\sqrt{m_\tau}\bigr)^2} = \frac{2}{3},
  \label{eq:koide}
\end{equation}
which holds to about one part in $10^5$ and has resisted explanation for forty years. The framework
produces it as follows. The three generations are the three off-diagonal octonionic entries of the
exceptional Jordan algebra $J_3(\mathbb{O})$ \autocite{baez2002octonions,duboisviolettetodorov2019exceptional},
cyclically permuted by a $\mathbb{Z}_3$ triality, so the square-root masses are the three triality
images of one quantity,
\begin{equation}
  \sqrt{m_k} = M\bigl(1 + \alpha\cos(\delta + \tfrac{2\pi k}{3})\bigr), \qquad k = 0,1,2,
  \label{eq:z3}
\end{equation}
with $M$ a scale, $\alpha$ the amplitude, and $\delta$ a phase. The amplitude is the Cayley--Dickson
doubling norm $\alpha = \sqrt2$ ($|1+i| = \sqrt2$), and this single value forces the Koide ratio.

\begin{proposition}[The Koide ratio from the doubling norm]\label{prop:koide}
For square-root masses of the form~\eqref{eq:z3}, the ratio~\eqref{eq:koide} equals
$(1 + \tfrac12\alpha^2)/3$, independently of $M$ and $\delta$. With $\alpha^2 = 2$ it equals $2/3$.
\end{proposition}

\begin{proof}
The three phases are spaced by $2\pi/3$, so $\sum_k \cos(\delta + 2\pi k/3) = 0$ and
$\sum_k \cos^2(\delta + 2\pi k/3) = \tfrac32$. Then $\sum_k\sqrt{m_k} = 3M$, so the denominator is
$9M^2$, while $\sum_k m_k = M^2\sum_k(1+\alpha\cos)^2 = M^2(3 + \tfrac32\alpha^2) = 3M^2(1+\tfrac12\alpha^2)$.
The ratio is $(1+\tfrac12\alpha^2)/3$, and $\alpha^2=2$ gives $2/3$.
\end{proof}

The Koide relation is thus derived, not fitted: it is the statement that the triality amplitude is
the doubling norm. The phase is $\delta = 2/9$ in radians, the value that the measured masses
require, matched to five significant figures and carrying a natural geometric reading as
$2/9 = \dim_{\mathbb{R}}\mathbb{C} \,/\, \dim_{\mathbb{R}}(\mathbb{O}\oplus\mathbb{R})$: the
preferred complex plane against the nine-dimensional ambient space of the Hopf base $S^8$ ---
exactly the (Bloch)$\,\oplus\,$(colour) space of Appendix~\ref{app:setup}. The reading is
suggestive, matched rather than derived. With $\alpha$ fixed by the theory, $\delta$ matched once,
and the single scale $M$ set by the heaviest mass, the other two follow:

\begin{center}
\begin{tabular}{lccc}
\hline
lepton & predicted (MeV) & observed (MeV) & error \\
\hline
$e$    & $0.5110$  & $0.5110$  & $-0.01\%$ \\
$\mu$  & $105.65$  & $105.66$  & $-0.01\%$ \\
$\tau$ & $1776.9$  & $1776.9$  & (scale) \\
\hline
\end{tabular}
\end{center}

This is the showcase: three masses from one dimensional input and no free dimensionless parameters,
the dimensionless content being the two theoretically fixed numbers $\alpha=\sqrt2$ and $\delta=2/9$.

We flag the one genuine gap squarely. Proposition~\ref{prop:koide} derives the $2/3$ relation from
$\alpha=\sqrt2$ and the $\mathbb{Z}_3$ ansatz, both of which the framework supplies; but the specific
phase $\delta = 2/9$ is at present a \emph{matched} value with a suggestive geometric interpretation,
not a derived one.

\begin{openproblem}[The Koide phase]\label{op:koide-delta}
Derive $\delta = 2/9$ from the holonomy of the octonionic Hopf connection. This is the single most
important open calculation in the mass sector. One route is already closed: the simplest
holonomy-squared model, $m_k \propto \sin^4(\theta_0 + 2\pi k/3)$, does not reproduce the spectrum
and is recorded as tried and failed.
\end{openproblem}

\subsection{Quarks: the colour correction}

Quarks do not satisfy the Koide relation --- the ratio~\eqref{eq:koide} is about $0.73$ for the
down-type and $0.85$ for the up-type quarks rather than $2/3$, a deviation that is invariant under
perturbative QCD running and therefore a bare structural fact, not a scale artefact. The framework
reads it as the effect of the colour debt that leptons, being colour singlets, do not carry. Two
modifications enter, one to the amplitude and one to the phase. Where the colourless leptons sit at
the bare doubling norm $\alpha^2 = 2$, the coloured quarks acquire an additional contribution that
scales with electric charge,
\begin{equation}
  \alpha^2 = 2 + 2\,|Q|^{3/2},
  \label{eq:quarkalpha}
\end{equation}
the coefficient being the doubling norm squared and the power $3/2$ fixed empirically to within one
per cent; the framework reads $|Q|^{3/2} = |Q|\cdot\sqrt{|Q|}$ as winding number times its
amplitude --- the Born-rule pairing appearing inside the fibre metric --- but this reading is a
gloss, not a derivation. The formula gives $\alpha^2 = 2.385$ for the down-type quarks
($|Q|=\tfrac13$) and $3.089$ for the up-type ($|Q|=\tfrac23$), against $2.389$ and $3.094$ extracted
from the measured sector ratios. The phase is \emph{diluted}: each sector's triality phase is
\begin{equation}
  \delta_{\mathrm{sector}} = \frac{2/9}{n},
  \qquad
  n \;=\; 1 + \bigl(T_3 + \tfrac12\bigr) +
  \begin{cases}1 & \text{colour triplet}\\[2pt] 0 & \text{colour singlet,}\end{cases}
  \label{eq:dilution}
\end{equation}
where $n$ counts the active fibre directions the channel engages --- the $S^1$ always, the $S^3$
for upper isospin, the $S^7$ for colour --- among which the misalignment $2/9$ is shared. Thus
$n = 1$ for the charged leptons ($\delta = 2/9$, as above), $n = 2$ for the down-type quarks
($\delta = \tfrac19$), and $n = 3$ for the up-type ($\delta = \tfrac{2}{27}$); triality shifts of
$2\pi/3$ relabel the generations without changing the spectrum.
Feeding~\eqref{eq:quarkalpha} and~\eqref{eq:dilution} into the triality form~\eqref{eq:z3}, with the
scale set by the heaviest quark in each sector, the lighter masses follow: $d$ at $4.62$~MeV
($-1.0\%$) and $s$ at $94.4$~MeV ($+1.1\%$) in the down sector; $c$ at $1271$~MeV ($+0.1\%$) in the
up sector. Eight of the nine charged fermions are reproduced to about one per cent.

The ninth is not. The up quark comes out at $2.78$~MeV against an observed $2.16 \pm 0.07$~MeV, an
error of some twenty-nine per cent and far outside the quoted uncertainty. This is a genuine failure
of the formula as it stands, not a rounding issue, and we record it as such: any claim that the
framework reproduces the quark spectrum must carry this exception openly. Two mitigating
observations follow, neither a resolution. The up-quark mass is never measured directly --- it is
extracted from lattice QCD with chiral perturbation theory, and determinations range over roughly
$1.9$--$2.5$~MeV, at the top of which range the discrepancy shrinks to about twelve per cent. And a
candidate mechanism is identified: at the first generation, and only there, the tripartite (GHZ) and
pairwise (W) contributions to the mass are of comparable size, and their interference shortens the
effective path --- consistent with the up quark also being the one fermion whose
$T_3 = +\tfrac12$ member is \emph{lighter} than its doublet partner, the competition between the
amplitude effect (larger $\alpha$ for up-type) and the dilution effect (smaller $\delta$) tipping
the other way at the lightest generation.

\begin{openproblem}[The up quark]\label{op:up-quark}
Give the generation-one GHZ--W interference mechanism a formal treatment and determine whether it
accounts for the $+29\%$ up-quark discrepancy, or find what does.
\end{openproblem}

There is a further structural thread that we flag rather than resolve. The quark--lepton mass
relations carry a factor of three of the Georgi--Jarlskog type \autocite{georgijarlskog1979new},
which the framework can read either group-theoretically, as a Casimir ratio of the colour structure,
or through the fitted charge dependence of~\eqref{eq:quarkalpha}. What the framework does \emph{not}
need is the $SU(5)$ scaffolding through which the Georgi--Jarlskog factor is conventionally obtained
--- consistent with the exclusion of $SU(5)$ and the absence of proton decay noted in
\S\ref{sec:charges}.

\begin{openproblem}[The colour correction]\label{op:colour-correction}
Derive the amplitude law~\eqref{eq:quarkalpha} --- in particular the power $3/2$ --- and the
phase-dilution rule~\eqref{eq:dilution} from the octonionic structure of the colour-triplet Higgs
coupling, and reconcile the Casimir and fitted accounts of the Georgi--Jarlskog factor.
\end{openproblem}

\subsection{Mass eigenstates versus weak eigenstates}

The basis that diagonalises the associator-debt --- the mass basis --- is not the basis in which the
weak interaction acts, which is the doublet basis of the $B$-qubit (\S\ref{sec:charges}). Their
mismatch is the Cabibbo--Kobayashi--Maskawa matrix. The framework ties the rotation to the same
octonionic geometry that fixes the masses: the lepton Koide phase $\delta = 2/9$ radians is about
$12.7^\circ$, strikingly close to the Cabibbo angle, which suggests that one geometry sets both the
mass hierarchy and the generation mixing. We mark that identification as suggestive rather than
derived.

\begin{openproblem}[The Cabibbo angle]\label{op:cabibbo}
Derive, rather than observe, the relation between the mass phase and the generation mixing: whether
the Cabibbo angle equals the undiluted phase $\delta = 2/9$, and what fixes the remaining CKM
structure.
\end{openproblem}

Its structural consequence, however, is exactly what the interaction picture needs: the mass
eigenstates are distinct from the weak eigenstates, so a charged-current vertex changes flavour while
the neutral currents, acting on the generation-blind structure, do not. That is the seed of the
no-tree-level flavour-changing-neutral-current result, which is delivered with the circuit picture of
interactions in \S\ref{sec:interactions}.

\subsection{What is settled, and what is owed}

Settled: mass is associator-debt, with the photon massless because it is debt-free; the Koide $2/3$
relation follows from the doubling norm $\alpha=\sqrt2$ and the $\mathbb{Z}_3$ triality of $J_3(\mathbb{O})$
(Proposition~\ref{prop:koide}); the charged-lepton spectrum is reproduced from one dimensional scale
and the two dimensionless numbers $\alpha=\sqrt2$, $\delta=2/9$, the electron and the muon each to
$0.01\%$; and the quark colour correction~\eqref{eq:quarkalpha}, with the phase-dilution
rule~\eqref{eq:dilution}, reproduces eight of nine charged fermion masses to about a per cent.

Owed, and recorded as the named problems above: the derivation of $\delta = 2/9$ from the Hopf
holonomy (Open Problem~\ref{op:koide-delta}, the key open calculation); the associator--geodesic
reconciliation (Open Problem~\ref{op:associator-geodesic}); the up-quark discrepancy
(Open Problem~\ref{op:up-quark}); the derivation of the colour correction and the dilution rule
(Open Problem~\ref{op:colour-correction}); the reduction of the sector scales
(Open Problem~\ref{op:sector-scales}); and the Cabibbo identification
(Open Problem~\ref{op:cabibbo}).

With the fermion masses in hand, two threads remain. The electroweak gauge bosons, whose masses use
the weak mixing angle $\sin^2\theta_W = \tfrac14$ already derived in \S\ref{sec:charges}, are taken
up in \S\ref{sec:boson-masses}; and the neutrino masses, which follow from the sterile seesaw of
\S\ref{sec:generations} and the dilution rule~\eqref{eq:dilution} with $n = 2$ --- a massless
lightest state, normal ordering, and a sum near $52$~meV --- are collected with the other
falsifiable predictions in \S\ref{sec:predictions}, where the one tension in the neutrino fit (the
first-to-second-generation splitting) is recorded as open.
